\subsection{几何对称、改进应力张量与守恒}

协变守恒律还须结合实际存在的 Killing 场才能产生守恒荷。下面先在正定几何中计算，再对照本章的 Lorentz 变分约定。

把对称张量视为纤维上齐次多项式 $T(x,z)=T_{a_1\cdots a_S}(x)z^{a_1}\cdots z^{a_S}$；引入 $z$ 仅为编码指标，所有等式独立于 $z$ 的数值。对迹自由张量，$\partial_z^2T=0$。定义协变对称梯度的未投影部分为
\[
 (\mathcal D T)(x,z)=z^a\nabla_aT(x,z),
\]
再通过无迹投影 $\operatorname{STF}$ 得到
\begin{equation}\label{eq:ghs-conformal-killing-tensor-operator}
 \mathcal K_S(T)=\operatorname{STF}(\nabla_{(a_0}T_{a_1\cdots a_S)}).
\end{equation}
$\mathcal K_S(T)=0$ 是一种符合几何的高阶 Killing 型方程；它是对背景对称性的约束，而不是任意张量自动满足的恒等式。

例如，在一阶 Killing 方程的计算中，
自旋一时，$\operatorname{STF}\nabla_{(a}\xi_{b)}=0$ 是共形 Killing 方程
\[
 \nabla_{(a}\xi_{b)}=\frac{1}{d}(\nabla_c\xi^c)g_{ab}.
\]
若右侧还为零，就得到严格 Killing 方程。由此可见，仅在存在这些特殊向量场时才可能利用它们构造守恒荷。

给定对称应力张量 $T^{ab}$，假定运动方程上 $\nabla_aT^{ab}=0$。若 $\xi$ 是向量场，则
\begin{align}
 \nabla_a(T^{ab}\xi_b)
 &=\underbrace{(\nabla_aT^{ab})\xi_b}_{0}
 +T^{ab}\nabla_a\xi_b\nonumber\\
 &=T^{ab}\nabla_{(a}\xi_{b)}.
 \label{eq:ghs-noether-killing-identity}
\end{align}
所以 Killing 向量产生守恒流；若 $\xi$ 仅共形 Killing，右侧变为 $T^a{}_a\nabla_b\xi^b/d$，需要应力张量在恰当意义下无迹才能守恒。这同时解释了为什么曲率下的 Weyl 异常会干扰经典共形荷。

令欧氏作用量
\[
 \mathscr I=\frac12\int_M\sqrt g\left[(\nabla\varphi)^2+(m^2+\xi\mathcal R)\varphi^2\right]\dd^dx.
\]
对度规取变分时，动能项给出 $\nabla_a\varphi\nabla_b\varphi-\frac12g_{ab}(\nabla\varphi)^2$；质量项贡献 $-\frac12g_{ab}m^2\varphi^2$。

曲率变分恒等式
\[
 \delta(\sqrt g\,\mathcal R)=\sqrt g\Bigl[G_{ab}\,\delta g^{ab}
 +(g_{ab}\Delta-\nabla_a\nabla_b)\delta g^{ab}\Bigr]
\]
在积分分部后产生改进项。采用 $T_{ab}=2\delta\mathscr I/(\sqrt g\,\delta g^{ab})$ 的符号约定，有
\[
 T_{ab}=\nabla_a\varphi\nabla_b\varphi
 -\frac12g_{ab}\bigl[(\nabla\varphi)^2+m^2\varphi^2\bigr]
 +\xi\bigl[G_{ab}\varphi^2+g_{ab}\Delta(\varphi^2)-\nabla_a\nabla_b(\varphi^2)\bigr].
\]
用运动方程 $(-\Delta+m^2+\xi\mathcal R)\varphi=0$、Bianchi 恒等式及 $[\Delta,\nabla_b]\varphi=\operatorname{Ric}_b{}^c\nabla_c\varphi$ 可直接检验 $\nabla^aT_{ab}=0$。

若把与欧氏作用量的物理压力惯例相反的应力张量定义用于其他章节，整体符号也须同步改变。

必须再次区分三个对象：局部张量自旋、全局等距群的角动量、动力学中守恒的高自旋流。它们在平直自由共形场论中高度重合，而在弯曲相互作用理论中通常分离。这种分离恰是一般化的实质内容。

高自旋场论中的抽象张量在初学时很难形成直观认识。其实最熟悉的自旋二对象已经存在于初等场论：能量动量张量。先在 $d>2$ 维欧氏平直空间考虑满足 $\Delta\varphi=0$ 的自由实标量场，以规范化作用量
$\mathscr I=\frac12\int(\partial\varphi)^2\dd^d x$ 为起点。平移不变性给出的对称张量是
\[
 T^{\rm can}_{\mu\nu}=\partial_\mu\varphi\partial_\nu\varphi
 -\frac12\delta_{\mu\nu}(\partial\varphi)^2.
\]
直接取散度，使用偏导交换与乘法法则，得到
\begin{align*}
 \partial^\mu T^{\rm can}_{\mu\nu}
 &=\underbrace{(\Delta\varphi)\partial_\nu\varphi}_{\text{运动方程项}}
 +\partial^\mu\varphi\,\partial_\mu\partial_\nu\varphi
 -\frac12\partial_\nu(\partial\varphi)^2\\
 &=(\Delta\varphi)\partial_\nu\varphi=0.
\end{align*}
但是它的迹一般不为零：$\delta^{\mu\nu}T^{\rm can}_{\mu\nu}=(1-d/2)(\partial\varphi)^2$。因此它并非直接适合共形表示论的无迹自旋二初级对象。

我们可在不改变守恒律的前提下加上一个恒等散度为零的改进项：
\begin{equation}\label{eq:ghs-bridge-improved-stress}
 T_{\mu\nu}=T^{\rm can}_{\mu\nu}
 +\xi_d(\delta_{\mu\nu}\Delta-\partial_\mu\partial_\nu)\varphi^2,
 \qquad \xi_d=\frac{d-2}{4(d-1)}.
\end{equation}
其散度为 $\partial^\mu(\delta_{\mu\nu}\Delta-\partial_\mu\partial_\nu)\varphi^2=0$，这是偏导可交换的恒等式，不必使用运动方程。迹则为
\begin{align*}
 T^\mu{}_{\mu}
 &=(1-d/2)(\partial\varphi)^2+\xi_d(d-1)\Delta\varphi^2\\
 &=(1-d/2+2\xi_d(d-1))(\partial\varphi)^2
 +2\xi_d(d-1)\varphi\Delta\varphi=0,
\end{align*}
其中用到 $\Delta\varphi^2=2(\partial\varphi)^2+2\varphi\Delta\varphi$。由此真正算到了两个不同的结论：守恒性来自平移对称，而无迹性需要特殊的改进系数与自由场运动方程。

对任何常向量 $b^\nu$，流 $J_\mu=T_{\mu\nu}b^\nu$ 满足 $\partial^\mu J_\mu=0$；若 $b^\nu$ 被共形 Killing 向量 $\xi^\nu(x)$ 取代，则
$\partial^\mu(T_{\mu\nu}\xi^\nu)=T_{\mu\nu}\partial^{(\mu}\xi^{\nu)}=(\partial\cdot\xi)T^\mu{}_{\mu}/d=0$。因此共形流的守恒从张量的两个可检验性质逐项推出。与此相比，一般曲率背景上的 $\nabla_\mu$ 不再可交换，平直改进项必须加入曲率修正；它们的来历可从前文协变导数的对易式追踪，而不能直接把 $\partial$ 替换成 $\nabla$。

这个自旋二算例也解释了高自旋课程中为何如此强调“短表示”：当守恒方程变为量子态条件时，散度后裔的范数为零；相互作用一般使该条件失效，相关算符的标度维数便离开幺正性界。后文的任意自旋计算将保持这种从运动方程和具体算符出发的方式。
